\problema{Critical Connections in a Network}{1192}{src/01-grafos/1192-critical-connections.cpp}{H}

Hay $n$ servidores numerados de $0$ a $n-1$, conectados por una lista de
\texttt{connections} (aristas \textbf{no dirigidas}). El grafo es
\emph{conexo}: siempre se puede llegar de cualquier servidor a cualquier
otro. Una conexión es \textbf{crítica} (un \emph{puente}) si, al quitarla, el
grafo se parte en dos o más componentes; es decir, es la única forma en que
un grupo de servidores se comunica con el resto. Debemos devolver
\textbf{todas} las conexiones críticas.

\paragraph{Intuición.} Piensa en cada conexión como un cable físico. Si un
cable forma parte de un ciclo (hay otra ruta alterna entre sus dos extremos),
quitarlo no aísla a nadie: la información puede rodear por el ciclo. Pero si
un cable es la \emph{única} ruta hacia cierta parte de la red —como una rama
que cuelga de un ciclo— quitarlo sí la desconecta. Ese cable es un puente.

\begin{center}
\ttfamily
\begin{tabular}{c}
\phantom{x}0\phantom{x}\\
/\ \ \textbackslash\\
1---2\\
$|$\\
3
\end{tabular}
\end{center}

\noindent Ejemplo con $n=4$, \texttt{connections} $=
[[0,1],[1,2],[2,0],[1,3]]$: el triángulo $0$-$1$-$2$ tiene tres rutas
distintas entre cualquier par de sus nodos, así que ninguna de sus tres
aristas es puente. Pero $1$-$3$ es la \emph{única} forma de llegar al
servidor $3$: quitarla lo aísla. La única conexión crítica es
$\{1,3\}$.

\paragraph{Solución: puentes de Tarjan.} Este es un patrón nuevo en el libro,
así que vale la pena explicarlo con calma. Hacemos un DFS desde cualquier
nodo (el grafo es conexo) y llevamos dos arreglos:

\begin{itemize}
  \item \texttt{disc[u]}: el ``número de orden'' en que el DFS
        \emph{descubrió} al nodo \texttt{u} por primera vez. Es un contador
        global que aumenta cada vez que visitamos un nodo nuevo (el primero
        visitado tiene \texttt{disc = 0}, el segundo \texttt{disc = 1}, \dots).
        Una vez asignado, \texttt{disc[u]} ya no cambia.
  \item \texttt{low[u]}: el \texttt{disc} \emph{más pequeño} al que se puede
        ``regresar'' desde el subárbol de \texttt{u}, usando aristas del
        árbol de DFS hacia descendientes y, como máximo, \textbf{una}
        arista de retroceso (\emph{back-edge}) que salte hacia un ancestro
        ya visitado. Este valor sí se actualiza (sólo puede bajar) conforme
        el DFS regresa de sus hijos.
\end{itemize}

Si al procesar la arista de árbol $(u, v)$ —con $v$ hijo de $u$— encontramos
una arista de retroceso desde el subárbol de $v$ hacia $u$ o algún ancestro
de $u$, entonces $v$ y su subárbol tienen una ``cuerda'' que los mantiene
conectados al resto del grafo aunque quitemos $(u,v)$: no es puente. Eso se
traduce en la condición clave:
$$(u,v) \text{ es puente} \iff \texttt{low[v]} > \texttt{disc[u]}$$
\texttt{low[v]} > \texttt{disc[u]} significa que, desde el subárbol de $v$,
es \emph{imposible} alcanzar a $u$ o a cualquier ancestro de $u$ salvo pasando
otra vez por la arista $(u,v)$: esa arista es la única puerta de entrada y
salida del subárbol, y por lo tanto es un puente.

En el ejemplo: al visitar $2$ (hijo de $1$) encontramos una arista de
retroceso hacia $0$ (ya visitado y no es su padre), así que
\texttt{low[2] = disc[0] = 0}; ese valor se propaga hacia arriba:
\texttt{low[1] = 0} y \texttt{low[0] = 0}. Ni $(0,1)$ ni $(1,2)$ son puente
porque $0 \not> \texttt{disc[0]}=0$ ni $0 \not> \texttt{disc[1]}=1$. En
cambio $3$ no tiene ninguna arista de retroceso (su único vecino es su
padre, $1$), así que \texttt{low[3]} se queda en su propio \texttt{disc = 3},
y $3 > \texttt{disc[1]} = 1$: la arista $(1,3)$ sí es puente.

\lstinputlisting{src/01-grafos/1192-critical-connections.cpp}

\complejidad{$O(V+E)$}{$O(V+E)$}

\paragraph{Notas de entrevista (Amazon).} Este problema evalúa si conoces
algoritmos de grafos más allá de BFS/DFS básico, así que vale la pena
explicar el porqué de \texttt{low[]}, no sólo memorizar la fórmula. Puntos
clave para mencionar: (1) al visitar una arista de retroceso, actualiza
\texttt{low[u]} con \texttt{disc[v]} (el momento de descubrimiento), \emph{no}
con \texttt{low[v]} — esa distinción es la que evita "atravesar" ciclos que
no existen. (2) Cuidado con las aristas paralelas hacia el padre: aquí se
evita re-visitar al padre comparando el nodo, pero si el grafo permitiera
múltiples conexiones idénticas entre los mismos dos servidores habría que
excluir sólo \emph{una} arista hacia el padre (por ejemplo, usando el índice
de la arista en vez del nodo padre); el enunciado de LeetCode garantiza que
no hay conexiones repetidas. (3) El orden en que se devuelven los puentes no
importa, pero normalizar cada arista como $(\min, \max)$ facilita comparar
resultados. Casos borde: $n=1$ sin conexiones (ningún puente posible), y un
árbol puro (sin ciclos), donde absolutamente todas las aristas son puentes.
